% Packages and notation used by the thesis content. Document presentation is
% controlled by the supplied MastersDoctoralThesis.cls.

\usepackage[utf8]{inputenc}
\usepackage[T1]{fontenc}
\usepackage{mathpazo}

\usepackage[backend=bibtex,style=authoryear,natbib=true,maxcitenames=2,mincitenames=1]{biblatex}
\addbibresource{refs.bib}
\usepackage[autostyle=true]{csquotes}

\usepackage{tikz}
\usepackage{tabularx}
\usepackage{booktabs}
\usepackage{multirow}
\usepackage{enumitem}
\usepackage{array}
\usepackage{amssymb}
\usepackage{amsmath}
\usepackage{amsthm}
\usepackage{graphicx}
\usepackage{xcolor}
\usepackage{makecell}
\usepackage{float}
\usepackage{subcaption}

\usepackage{etoolbox}

\makeatletter
\patchcmd{\l@chapter}{\vskip 1.0em}{\vskip 0.3em}{}{}
\renewcommand*\l@figure{\@dottedtocline{1}{0em}{3em}}
\renewcommand*\l@table{\@dottedtocline{1}{0em}{3em}}
\setlength{\@fptop}{0pt}
\setlength{\@fpbot}{0pt plus 1fil}
\makeatother

% ---------------------------------------------------------------------------
% Notation macros (see docs/paper-outline.md Ch.3 for the formal definitions)
% ---------------------------------------------------------------------------
% Paper's own reported number
\newcommand{\Pref}{P}
% Agent signal + agent-read paper configuration (paper-faithful baseline)
\newcommand{\Aref}{A}
% Agent signal + C&Z's actual configuration (config-swapped hybrid track)
\newcommand{\Acz}{A_{cz}}
% Agent signal + HXZ-style standardized configuration (config-swapped hybrid track)
\newcommand{\Ahxz}{A_{hxz}}
% C&Z's own measured/reported result (from SignalDoc.csv / CZReferenceProfile)
\newcommand{\CZ}{CZ}
% HXZ's own reported result for this factor (via hxz_bridge.compute_hxz_reported)
\newcommand{\HXZ}{HXZ}
% Agent-extracted config / C&Z's actual config / HXZ standardized config
\newcommand{\Cagent}{C_{\text{agent}}}
\newcommand{\Ccz}{C_{\text{cz}}}
\newcommand{\Cstd}{C_{\text{std}}}
% Agent signal series
\newcommand{\SA}{S_A}
% t-statistic shorthand
\newcommand{\tstat}{t}
% Implementation space admitted by a paper P
\newcommand{\ImplSpace}[1]{\mathcal{I}(#1)}

\theoremstyle{definition}
\newtheorem{definition}{Definition}
\newtheorem{assumption}{Assumption}
